\begin{INSRContentUnit}[
  id = introduction,
  title = {Introduction},
  outputs = {paper,position-paper,report,slides,handout}
]

\INSRFullText{
Complex post-traumatic stress disorder is associated not only with core post-traumatic symptoms, but also with disturbances in affect regulation, self-organization, interpersonal functioning, and state-dependent access to cognitive and emotional resources. Although established trauma-focused treatments can be effective, their clinical delivery frequently depends on moment-to-moment decisions about pacing, stabilization, exposure readiness, dissociation, mentalization, autonomic activation, and therapeutic alliance.

These decisions are commonly guided by clinical expertise, yet they are not always represented within an explicit and testable decision architecture. INSR addresses this gap by proposing that interventions should be selected and sequenced according to the patient's current regulatory, cognitive, affective, relational, and dissociative state rather than solely according to diagnosis or a fixed treatment phase.

The purpose of this position paper is to define the conceptual structure of INSR, distinguish it from a conventional multimodal treatment package, and identify the elements that require operationalization and empirical validation.
}

\INSRSummary{
INSR formalizes clinical decisions about pacing, stabilization, processing readiness, and safety in complex trauma treatment.
}

\INSRKeyMessage{
INSR is a decision architecture for determining when an intervention is appropriate, not merely a collection of therapeutic techniques.
}

\end{INSRContentUnit}
